\chapter{Diskussion}
\label{ch:diskussion}
Nun werden die Ergebnisse der Evaluation eingeordnet. Danach folgen die Einschränkungen der Arbeit und Handlungsempfehlungen für den praktischen Einsatz lokaler Sprachmodelle in mobilen BCI-Anwendungen.

\section{Einordnung der Ergebnisse}
\label{sec:einordnung}
Die verwendeten Gemma-Modelle unterscheiden sich in ihren Stärken. Gemma 3 1B ist aufgrund seiner geringen Größe am effizientesten und hat den geringsten Speicher- und Akkubedarf sowie die höchste Geschwindigkeit. Das liegt daran, dass die Modellgröße einen Einfluss auf den Speicher- und Rechenaufwand hat (\autoref{subsec:ondevice}). Gemma 2 2B liegt innerhalb der beiden Gemma-Modelle an erster Stelle, wenn man allein die Ausgabequalität betrachtet.
\vspace{\baselineskip}

\noindent Qwen 2.5 1.5B zeigt, dass die Größe allein die Qualität nicht bestimmt. Obwohl es sich von der Parameterzahl her in der Mitte der drei Modelle einordnet, liefert es die schlechtesten Ausgaben. Die technischen Bedingungen waren für alle Modelle gleich (\autoref{ch:methodik}).
Qwen 2.5 1.5B hat die höchste Inferenzdauer und verbraucht am meisten Akku, obwohl es kleiner als Gemma 2 2B ist. Das Modell pendelt sich aber bei dem Rest, bis auf die Ladezeit im Artefakt- und Vergleichsszenario, bei der es mit 0,27\,s am schnellsten ist, in der Mitte ein. Der Grund für die lange Inferenzdauer liegt in der Länge der Antworten. Je mehr Text ein Modell erzeugt, desto länger dauert die Berechnung und desto mehr Energie verbraucht sie. Qwen 2.5 1.5B erzeugt in den Messungen deutlich mehr Tokens (vgl.\ \autoref{tab:app-rohdaten}). Wenn man den Durchsatz von Qwen 2.5 1.5B betrachtet, liegt er mit 11,31 bis 14,06 Tokens pro Sekunde in einem deutlich schnelleren Bereich als Gemma 2 2B mit 2,87 bis 5,06 (\autoref{tab:eval-latenz}). Qwen 2.5 1.5B erzeugt seine Tokens also nicht langsamer, sondern nur mehr davon. Die Ausführlichkeit ist gleichzeitig ein Qualitätsproblem, weil es sich schlecht an die Vorgaben hält. Trotz der klaren Formatvorgaben im Prompt fällt die Antwort zu lang aus. Die Ausgabe ist so nicht nur inhaltlich am schwächsten, sondern kostet auch am meisten Zeit und Energie. Beides tritt bei Qwen gemeinsam auf, möglicherweise durch das Nichteinhalten der Längenvorgaben oder eine ungeeignete Promptwahl. Ein fester Zusammenhang lässt sich jedoch aus den Messungen nicht ableiten.
\vspace{\baselineskip}

\noindent Die \ac{TTFT} ist im Vergleichsszenario bei allen drei Modellen deutlich länger als bei Report und Artefakt-Erklärung. Das liegt an der längeren Prompteingabe, bei der gleich zwei Merkmalstexte (einer je Session) verarbeitet werden. Je mehr Text ein Modell liest, desto länger dauert es bei diesem Modell in der Regel, bis die Antwort generiert wird. Dies bestätigt sich in den Eingabe-Tokenzahlen (\autoref{tab:eval-eingabetokens}), die im Vergleichsszenario am größten ausfallen. Wichtig zu erwähnen ist, dass Qwen 2.5 1.5B den Prompt in allen Szenarien in die meisten Tokens zerlegt. Der Zerlegungsunterschied zwischen den Modellen kann durch das kleinere Vokabular von Qwen 2.5 1.5B (\autoref{subsec:messbedingungen}) entstehen, ebenso können die abweichenden Zerlegungsverfahren oder die Struktur des Prompts die Ursache sein. Welcher dieser Einflüsse überwiegt, bleibt offen, da es aus den Daten nicht sicher abgeleitet werden kann. Da die Tokenanzahl der Eingabe gegenüber Gemma 3 1B nur wenige Prozent (4 bis 8) höher ist, lässt sich die im Verhältnis klar erhöhte \ac{TTFT} von Qwen 2.5 1.5B nicht mit der Länge der Eingabe erklären. Gemma 2 2B erreicht mit nahezu gleich vielen Eingabe-Tokens wie Gemma 3 1B sogar die höchste \ac{TTFT}. Damit begründet sich der Unterschied zwischen den Modellen wahrscheinlich durch den erhöhten Rechenaufwand (\autoref{subsec:ondevice}). Innerhalb eines Modells können die Eingabe-Tokens einen Einfluss auf die \ac{TTFT} haben, zwischen den Modellen erklären sie die Unterschiede jedoch nicht. Eine Ausnahme zeigt sich bei Qwen 2.5 1.5B im Artefakt-Szenario, wo die \ac{TTFT} trotz mehr Eingabe-Tokens leicht unter dem Report-Szenario liegt. Der Unterschied ist mit 0,05\,s aber sehr gering und liegt innerhalb der Standardabweichung (\autoref{tab:eval-latenz}).
\vspace{\baselineskip}

\noindent Die Tokens pro Sekunde werden direkt von der längeren Vorlaufzeit beeinflusst. Die Metrik ergibt sich aus den erzeugten Tokens geteilt durch die gesamte Inferenzdauer (\autoref{subsec:metriken}), bei der die \ac{TTFT} einen großen Anteil ausmachen kann. So entfallen im Vergleichsszenario bei Gemma 3 1B rund 59 Prozent und bei Gemma 2 2B rund 61 Prozent der Inferenzdauer auf die Phase vor dem ersten Token, im Report-Szenario sind es 29 beziehungsweise 43 Prozent. Da bei Qwen 2.5 1.5B der Anteil aufgrund der längeren Ausgaben im Report- und Vergleichsszenario bei nur 16 beziehungsweise 33 Prozent liegt, fällt dort auch der Durchsatz am wenigsten ab. \autoref{tab:eval-latenz} zeigt, dass sich der Durchsatz vom Report- zum Vergleichsszenario bei dem Modell Gemma 3 1B von 21,65 auf 11,85 Tokens pro Sekunde fast halbiert. So auch bei Gemma 2 2B, bei dem der Durchsatz von 5,06 auf 2,87 sinkt. Bei Qwen 2.5 1.5B sinkt er nur von 14,06 auf 11,31 Tokens pro Sekunde. Dass der Rückgang an der Vorlaufzeit liegt, zeigt sich, wenn man sie herausrechnet. Dann bleibt die reine Erzeugungsrate weitgehend gleich, bei Gemma 3 1B 30,5 gegenüber 29,0 und bei Qwen 2.5 1.5B 16,7 gegenüber 16,9 Tokens pro Sekunde. Spürbar sinkt sie nur bei Gemma 2 2B, von 8,8 auf 7,4, da sich die Werte relativ zueinander um rund 16 Prozent unterscheiden. Der Einbruch geht damit vor allem auf die längere Verarbeitung der Eingabe zurück. Die Ladezeit liegt im Bereich von 0,27 bis 0,87\,s bei allen Modell-Szenario-Kombinationen und fällt daher gegenüber der bis zu 42\,s langen Inferenzdauer kaum ins Gewicht (\autoref{tab:eval-latenz}).
\vspace{\baselineskip}

\noindent Durch den Verbrauch in \ac{PSS} kann man den von der Anwendung benutzten \ac{RAM} ins Verhältnis zum gesamten \ac{RAM} von 8 GB setzen. Gemma 2 2B belegt mit ca.\ 3,7 GB rund 45 Prozent der verfügbaren 8 GB, Gemma 3 1B mit ca.\ 1,5 GB rund 18 Prozent und Qwen 2.5 1.5B mit ca.\ 2,0 GB rund 24 Prozent. Bei dem 2B-Modell wären es auf einem Gerät mit 4 GB Arbeitsspeicher ca.\ 90 Prozent und damit an der Grenze. Da der \ac{PSS}-Wert eine obere Schranke darstellt (\autoref{subsec:metriken}), sind dies Höchstwerte. Zu beachten ist, dass einer Anwendung nie der gesamte Arbeitsspeicher netto zur Verfügung steht, da das System den Speicher unter Systemprozessen und Anwendungen aufteilt und bei Knappheit Prozesse beendet \autocite{android_memory}. Daraus lässt sich ableiten, dass 4 GB für das 2B-Modell sehr wahrscheinlich zu klein wären. Wahrscheinlich sind 6 GB oder 8 GB das Mindestmaß, das für das 2B-Modell benötigt wird. Geräte mit mehr \ac{RAM} können je nach Kompatibilität, Hardware und Android-Version größere Modelle verwenden.
\vspace{\baselineskip}

\noindent Mögliche Gründe für das schlechte Einhalten der Vorgaben bei Qwen lassen sich aus den Messungen allein nicht klären. Das Kontextfenster ist ausgeschlossen, da selbst im Vergleichsszenario mit 770 Eingabe- und 474 Ausgabe-Tokens die Grenze von 1280 Tokens nicht erreicht wird (\autoref{tab:app-rohdaten}). Wahrscheinlich liegt es an der Feinjustierung des Prompts, die speziell an das Modell Gemma 2 2B angepasst wurde (\autoref{sec:einschraenkungen}). Denkbar ist zudem, dass sich die Modelle darin unterscheiden, wie strikt sie Längen- und Formatvorgaben umsetzen.
\vspace{\baselineskip}

\noindent Hinzu kommt, dass Qwen 2.5 1.5B einen größeren Nachteil hat als die Gemma-Modelle. Gemma 3 1B kommt aus derselben Modellfamilie wie der Vorgänger und hat dadurch wahrscheinlich einen größeren Mehrwert durch die Prompt-Feinjustierung des Vorgängers. Qwen kommt von einem anderen Hersteller mit einem eigenen Tokenizer (\autoref{sec:qwen}), wodurch eine negative Beeinflussung nicht ausgeschlossen ist. Es kann auch sein, dass Qwen gegenüber den verwendeten Gemma-Modellen eine 8-Bit-Quantisierung und den Standardwert der Temperatur anders verarbeitet (\autoref{subsec:ondevice}, \autoref{sec:modellauswahl}). Die schwächeren Ergebnisse von Qwen sind daher nicht allein als geringere Fähigkeit zu lesen, sondern zum Teil als Folge eines Prompts, der nicht auf dieses Modell abgestimmt ist, und möglicherweise der Quantisierung oder der Temperatur.
\vspace{\baselineskip}

\noindent Die Ergebnisse zeigen somit, dass sich von der Modellgröße nicht unmittelbar auf die Ausgabequalität schließen lässt. Verschiedene Hersteller unterscheiden sich in Aufbau und Training der Modelle \autocite{gemma2_2024, qwen2025}, sodass die Größe als Anhaltspunkt nicht verlässlich genug ist. Bei der Stabilität unterscheiden sich die Modelle hingegen nicht. Alle Kombinationen laufen fehlerfrei durch. Die Feststellung in \autoref{sec:eval-latenz}, dass die Ausgabe immer identisch ist, liegt am Zufallsstartwert der Schnittstelle (\autoref{sec:modellauswahl}). Nur Qwen 2.5 1.5B weist einen erhöhten Streuungswert im Vergleichsszenario auf, was mit seiner langen Laufzeit unter Dauerlast zusammenhängen kann. Für die Forschungsfrage ist das Wichtigste, dass die Modelle einwandfrei und stabil ihre Aufgaben erfüllen und einen Mehrwert bieten. Solange die Latenz in einem vertretbaren Rahmen liegt, ist die Ausgabequalität das wichtigste Kriterium. Ob das hier bei den genutzten Modellen zutrifft, lässt sich nur vermuten.

\section{Einschränkungen}
\label{sec:einschraenkungen}
Während der gesamten Arbeit gibt es Einschränkungen, die die Ergebnisse beeinflussen können. Dazu gehören unter anderem Limitierungen durch Hardware, Designentscheidungen oder Entscheidungen bezogen auf die Messungen.
Die Arbeit basiert auf nur einem Testgerät, dem Samsung Galaxy S23. Es ist nicht sicher, wie die Evaluation auf anderen Hardwarekonfigurationen ausfallen würde. Das Testgerät schränkt zudem die Auswahl der Modelle ein. Aufgrund seiner begrenzten Unterstützung ist es nicht möglich, systemintegrierte Modelle wie Gemini Nano einzubeziehen \autocite{mlkit_genai_devices}. Zwar unterstützt das Testgerät Galaxy AI, dabei handelt es sich jedoch um vordefinierte Endnutzerfunktionen \autocite{samsung_galaxyai} und nicht um eine Schnittstelle, über die eigene Prompts an ein Sprachmodell übergeben werden könnten. Für den hier untersuchten Anwendungsfall ist es deshalb nicht einsetzbar. Auch bei den frei deploybaren Modellen ist die Auswahl begrenzt. Durch die selbst getroffene Wahl von MediaPipe als Laufzeitumgebung entfallen viele neuere Modelle, da die Schnittstelle nur eine begrenzte Auswahl unterstützt \autocite{mediapipe_llm} oder die Modelle für mobile Geräte gar nicht optimiert sind \autocite{edgellm_review}. Zudem wird für einen technisch einheitlichen Vergleich darauf geachtet, dass alle mit einer Kontextgröße von 1280, einer 8-Bit-Quantisierung und dem .task-Format laufen, was zur Folge hat, dass noch mehr Modelle ausgeschlossen werden.
\vspace{\baselineskip}

\noindent Die gesamte Promptfeinjustierung basiert auf dem Modell Gemma 2 2B, wodurch das Setup auf dieses Modell zugeschnitten ist. Davon betroffen ist sehr wahrscheinlich die Ausgabequalität sowie aufgrund der Länge der Ausgabe die Inferenzdauer, der Durchsatz und der Akkuverbrauch. Da die betroffenen Werte als Ergebnis unter günstigen Bedingungen entstanden sind, stellt sich die Frage, ob Gemma 2 2B in den Vergleich gezogen werden soll. Da Ladezeit, \ac{TTFT} und Speicherbedarf davon unberührt bleiben, wird das Modell im Vergleich belassen. Die betroffenen Werte sind jedoch entsprechend einzuordnen. Der Zuschnitt wirkt zudem nicht nur auf Gemma 2 2B. Da Gemma 3 1B derselben Modellfamilie angehört, dürfte es von der Feinjustierung teilweise mitprofitieren, während Qwen 2.5 1.5B als Modell eines anderen Herstellers keinen solchen Vorteil hat.\vspace{\baselineskip}

\noindent Beim Messaufbau ist die Reihenfolge der Messungen bewusst festgelegt (\autoref{sec:messprotokoll}). Trotz einer Abkühlphase lässt sich nicht ausschließen, dass die zuletzt gemessenen Modelle durch das wärmere Gerät leicht benachteiligt werden (\autoref{subsec:ondevice}). Dasselbe gilt für den Ladezustand, der über die Messreihe hinweg sinkt. Die Streuung der Zeiten geht deshalb auf die Messumgebung zurück. Der verbleibende Einfluss bleibt offen. Alle Messungen erfolgen zudem in deutscher Sprache, daher bleibt es offen, wie sich die Ausgabequalität und Tokenanzahl in englischer Sprache verhalten. Schließlich sind die drei Szenarien in ihrer Aufgabenart identisch, nämlich in der textuellen Beschreibung von \ac{EEG}-Sessions (\autoref{sec:szenarien}). Der Unterschied liegt im Anspruch und nicht in der Art der Aufgabe. Daher bleibt es offen, wie die Modelle bei anderen Aufgaben wie der Klassifikation abschneiden würden.
\vspace{\baselineskip}

\noindent Wie in \autoref{subsec:aufbereitung} beschrieben, wächst der Aufwand der Merkmalsextraktion linear mit der Länge der Session. Die Inferenzdauer ist davon nicht beeinflusst, da unabhängig von der Länge immer vier Segmente entstehen. Bei längeren Sessions von vielen Minuten könnte sich ihr Einfluss deutlich erhöhen und die Extraktion stärker ins Gewicht fallen.

\section{Handlungsempfehlungen}
\label{sec:empfehlungen}
In diesem Abschnitt werden Handlungen auf Basis der Erkenntnisse dieser Arbeit empfohlen. Diese leiten sich aus den Ergebnissen der Evaluation und aus der Umsetzung ab. Die folgenden Empfehlungen gelten für den in dieser Arbeit untersuchten Rahmen, also für die beschriebene Implementierung mit der in \autoref{sec:impl-prompt} festgelegten Prompt-Struktur und der gewählten Konfiguration, für die drei Anwendungsszenarien aus \autoref{sec:szenarien} und für das Testgerät aus \autoref{subsec:messbedingungen}. Auf andere Geräte, Prompts oder Aufgaben lassen sie sich nicht ohne Weiteres übertragen.

\subsection{Allgemeine Empfehlungen}
Die Prompts sollten im besten Fall kurz und klar gehalten werden, damit sie in den Kontext passen. Die Länge beeinflusst ebenso die Zeit bis zum ersten Token (\autoref{sec:einordnung}).
Das Modell sollte die Merkmale in einer klar verständlichen Form erhalten, ohne selbst überfordert zu sein, also keine Rohdaten, sondern aggregierte Merkmale.
Darüber hinaus sollte man je nach Szenario das passende Modell auswählen, da die Ausgabequalität je nach Komplexität der Aufgabe variiert.
Da die Modelle jeweils 1,5\,GB, 2,0\,GB und 3,7\,GB im Arbeitsspeicher belegen (\autoref{tab:eval-ressourcen}), sollten sie erst nach der Prüfung der Eingabe geladen werden. Nach der Ausgabe können sie schließlich wieder freigegeben werden.
Die Inferenz der Modelle sollte auf einen Hintergrund-Thread verlagert werden, da eine Anfrage bis zu 42\,s dauern kann und die Oberfläche einfrieren würde (\autoref{tab:eval-latenz}).
Das Modell selbst sollte hinter einer eigenen Schnittstelle gekapselt werden. In dieser Arbeit ließen sich die drei Modelle allein über den Modellpfad austauschen (\autoref{sec:architektur}), sodass ein Wechsel des Modells oder der Laufzeitumgebung die übrige Anwendung nicht betrifft.
Für die Benutzerfreundlichkeit ist es auch sinnvoll, eine wie hier verwendete Streaming-Mechanik einzubinden, um die gefühlte Wartezeit zu verringern. Die Zeit wird dadurch nicht kürzer, aber vermutlich so empfunden. Dabei sollte bei der Ausgabe das Format nach Bedarf angepasst werden.
Bei Qwen 2.5 1.5B ist es beispielsweise sinnvoll, die Antwort zu kürzen und formatierungstechnisch zu vereinfachen. Selbst das größte Modell Gemma 2 2B hat keine fehlerfreien Ausgaben geliefert, weshalb man sich in Bereichen wie der Medizin nicht allein auf solche Modelle im Rahmen dieser Arbeit verlassen sollte.

\subsection{Modellempfehlung}
Wenn der Fokus auf der Qualität liegen soll, geht die Empfehlung an das Modell Gemma 2 2B. Es hält sich am besten an die Vorgaben und bietet die beste Ausgabequalität. Dafür verbraucht es am meisten Speicher und ist am langsamsten, mit einer geringen Token-Rate und der längsten \ac{TTFT}. Aufgrund des Speicherbedarfs von ca.\ 3,7\,GB für den gesamten Prozess der Anwendung empfiehlt es sich, ein Gerät mit mindestens 6 bis 8\,GB Arbeitsspeicher zu nutzen (\autoref{sec:einordnung}). Mit Betriebssystem, einem Puffer und möglicherweise weiteren Anwendungen kann der Speicherbedarf ansteigen.
\vspace{\baselineskip}

\noindent Wer hingegen auf Geschwindigkeit, Speicher- und Akkuverbrauch setzen möchte, ist mit dem kleinsten Modell Gemma 3 1B am besten aufgehoben. Es ist für Geräte mit wenig Arbeitsspeicher die beste Wahl, da das Modell in diesem Punkt die geringsten Anforderungen stellt.
\vspace{\baselineskip}

\noindent Qwen 2.5 1.5B ist in diesem Rahmen nicht zu empfehlen. Es schneidet in puncto Qualität und Akkuverbrauch trotz seiner Größe am schlechtesten ab. Die Geschwindigkeit und der Speicherbedarf pendeln sich, mit Ausnahme der Inferenzdauer, wie erwartet in der Mitte ein, können jedoch nicht die mangelnde Ausgabequalität rechtfertigen. Dabei ist zu berücksichtigen, dass der Prompt an Gemma 2 2B ausgerichtet wurde und Qwen dadurch benachteiligt ist (\autoref{sec:einschraenkungen}).
\vspace{\baselineskip}

\noindent Je nach Einsatzzweck haben die beiden Gemma-Modelle in dieser Untersuchung Vorteile. Für schnelle Antworten mit einer geringen \ac{TTFT} sowie einer niedrigen Inferenzdauer liegt Gemma 3 1B vorne. Wenn hingegen Qualität wichtig und Zeit gegeben ist, beispielsweise als Hintergrundprozess ohne nötige sofortige Rückmeldung, ist Gemma 2 2B die beste Wahl.